% Part: quantified-modal-logic
% Chapter: semantics
% Section: classical-rules

\documentclass[../../../include/open-logic-section]{subfiles}

\begin{document}

\olfileid{qml}{nd}{cla}

\olsection{Classical Quantified Modal Logic}

In Natural Deduction !!{derivation}s for classical 
quantified modal logic (\Log{CQML}), you are allowed to use the following rules, detailed
in the relevant sections of this textbook:

\begin{enumerate}
  \item Propositional Logic. All rules of propositional logic.
  \item Classical Quantification and Identity. All rules of
  first-order logic.
  \item Existence. The definition of the existence predicate
  $\lfrexists$ (see below).
  \item Modality. All rules of modal logics. Depending on how the
    accessibility 
    relation of our models are constrained:
    \begin{enumerate}
      \item By default, \Log{CQML} uses modal rules from the basic
      modal system \Log{K}.
      \item Stronger systems are \Log{CQML+X}, for instance
      \Log{CQML+T} is classical quantified modal logic with modal rules
      from system \Log{T}. Its models are fixed-domain models where
      the accessibility relation is reflexive.
    \end{enumerate}
\end{enumerate}

\begin{defish}
    \emph{Definition of $\lfrexists$}

    \begin{center}        
    \smallskip
    \begin{tabular}{cc}
    \AxiomC{$\lexists[v][\eq[c][v]]$}
    \RightLabel{Def~$\lfrexists$}
    \UnaryInfC{$\lfrexists c$}
    \DisplayProof
    &
    \AxiomC{$\lfrexists c$}
    \RightLabel{Def~$\lfrexists$}
    \UnaryInfC{$\lexists[v][\eq[c][v]]$}
    \DisplayProof
    \end{tabular}
    \end{center}

\end{defish}

\end{document}